\section{Discussion and Social Impact}
\label{sec:discussion}

PACT targets a concrete governance decision. A bank, hospital, or
municipal office choosing an agent for a regulated workflow has no instrument that
answers ``which model holds this rule under a rushed manager?'' A per-domain,
floor-sensitive, multi-turn propensity profile gives a deployer something to act on:
filter the leaderboard to their own domain, read the worst-quartile floor rather
than the flattering mean, check whether a guardrail actually moves the candidate
(steerability), and see whether a violating model at least admits it (honesty). Two
findings are already interventions: rephrasing a rule imperatively raises
compliance, and removing quantitative penalty language avoids the ``fine is a
price'' paradox for prone models. The frozen item set, with its pre-registered hash,
private holdout, and canary, also lets the same benchmark rerun on each model
\emph{version}; no one currently tracks whether a vendor's silent update moved its
urgency floor, and a benchmark that catches that becomes monitoring infrastructure.
Because the release is per-trial and every aggregator ships, the headline is a
recommendation with an argument, and a deployer with a different cost model for
violations versus escalations can re-rank.

\paragraph{Threats to validity.} \emph{Evaluation awareness} would undermine any
propensity benchmark; we enforce naturalism at generation and report a dedicated
awareness probe, measuring the threat rather than assuming it away. \emph{Stated
versus enacted compliance}: recommendation-level compliance upper-bounds what a
tool-using agent actuates, and an actuated variant is future work, but \emph{Moffatt
v.\ Air Canada} established that a chatbot's statement is itself a liability
surface~\citep{moffatt2024aircanada}. \emph{Judge circularity} is met by the
deterministic first tier, the human gold set, and the judge-swap $\tau$-gate.
\emph{Generated items} could favor their author's family, which the dual-guard
review and the computable generator-versus-evaluatee test address. \emph{Panel
size}: model-level analyses are indicative at twelve models, so the item-level
statistics over thousands of trials are load-bearing, and every new model
strengthens them.

\section{Conclusion}
\label{sec:conclusion}

PACT measures whether enterprise agents keep an embedded rule when benign
incentive conflict, institutional pressure, and multi-turn pushback make violation
attractive. Its contribution is methodological and empirical: a template comparable
across twelve regulated domains, a theory-grounded pressure battery, a two-arm
multi-turn protocol yielding six de-correlated axes, an auditable and
contamination-resistant generation pipeline, and a metric suite validated on a
72{,}000-trial precursor and calibrated to resist degenerate policies. The failure
class has already produced tribunal damages, a certified class action, and regulator
orders, and the instrument reruns cheaply as models change. We release the frozen
items, per-trial outcomes, reviewer logs, and every aggregator so the benchmark can
be re-analyzed, extended through its promotion gate, and used as what regulated
deployers lack: a standard way to ask an agent, before trusting it with a rule,
whether it will keep the rule when keeping it costs something.

\section*{Ethical Statement}
This is a defensive evaluation: it measures whether AI agents follow regulations and
policies, to help deployers and regulators identify models that fail to. All
scenarios are synthetic; no real personal data or confidential business information
is used, and the litigated incidents referenced are public record. Characterizing
how models can be pushed into non-compliance could in principle inform an actor
seeking to elicit it, but the pressures we catalog are ordinary workplace situations
already ubiquitous in deployment, not novel attack techniques, and the value of
letting deployers detect non-compliant agents before trusting them with regulated
decisions outweighs this risk. We release with a private holdout to preserve the
benchmark's integrity as a governance instrument.
